\section{Abstract}\label{abstract}

To inform public health interventions, researchers have developed models to forecast opioid-related overdose mortality. However, these efforts often have limited overlap in the models and datasets employed, presenting challenges to assessing progress in this field. Furthermore, common error-based performance metrics, such as root mean squared error, are not directly suitable to assess a key modeling purpose: the identification of priority areas for public health interventions. We recommend a new intervention-aware performance metric and establish a set of baseline models with competitive performance. \hl{To show how model and metric choice vary across locations, we explore two distinct geographies: Cook County, Illinois and the state of Massachusetts. We introduce a new, intervention-aware evaluation metric: the Percentage of Best Possible Reach (\%BPR).} The top-performing models based on error-based metrics recommend fixed-budget interventions in areas that do not always reach the most possible overdose events. In Massachusetts the top models, as ranked by our proposed \%BPR metric, could have reached 18 additional fatal overdoses per year in our 2020-2021 test period compared to models favored by error-based metrics, assuming the ability to intervene in 100 census tracts out of the 1620 in Massachusetts. We release open code and data for others to build upon.

Repository for code and data: \href{https://github.com/tufts-ml/opioid-overdose-models}{\ul{https://github.com/tufts-ml/opioid-overdose-models}}
